% Annale FICTIVE (aucun contenu réel) : examen français, exercices numérotés en romain, barème entre parenthèses.
\documentclass[11pt,a4paper]{article}
\usepackage[T1]{fontenc}
\usepackage[utf8]{inputenc}
\usepackage{lmodern}
\usepackage[margin=2cm]{geometry}
\usepackage{fancyhdr}
\usepackage{lastpage}
\pagestyle{fancy}\fancyhf{}
\fancyhead[L]{\small INV 300 -- Algèbre inventée}
\fancyhead[R]{\small Examen terminal}
\fancyfoot[R]{\small \thepage/\pageref{LastPage}}
\begin{document}
\thispagestyle{empty}
\begin{center}
{\large Université Fictive de Testville}\\[8pt]
{\Large\textbf{INV 300 -- Algèbre inventée}}\\[4pt]
Examen terminal -- Durée : 2 heures
\end{center}
\vspace{8pt}
\noindent Nom : \makebox[0.4\textwidth]{\dotfill}\hfill Prénom : \makebox[0.3\textwidth]{\dotfill}\par\vspace{6pt}
\noindent Numéro d'étudiant : \makebox[0.3\textwidth]{\dotfill}\par\vspace{10pt}
\noindent\textbf{Consignes}
\begin{itemize}\setlength{\itemsep}{1pt}
\item Aucun document n'est autorisé.
\item Les calculatrices sont interdites.
\item La qualité de la rédaction sera prise en compte.
\end{itemize}
\vspace{10pt}
\noindent\textbf{Exercice I} (6 points)\par\smallskip
Soit $u_0 = 1$ et $u_{n+1} = 3u_n$. Calculer $u_4$.\par\vspace{3cm}
\noindent\textbf{Exercice II} (7 points)\par\smallskip
Montrer que la somme de deux nombres pairs inventés est paire.\par\vspace{3cm}
\clearpage
\noindent\textbf{Exercice III} (7 points)\par\smallskip
Résoudre $2x + 3 = 11$ dans un corps fictif de caractéristique nulle.\par\vspace{3cm}
\end{document}
